%% ============================================================
%% School Report Template (Japanese / LuaLaTeX)
%% Compile with: lualatex main.tex  or  latexmk main.tex
%%
%% \documentclass[options]{class} is always the first command.
%% '12pt' sets the base font size; 'a4paper' sets the paper size.
%% 'ltjsarticle' is the LuaLaTeX Japanese article document class.
%% ============================================================
\documentclass[12pt,a4paper]{ltjsarticle}

%% ------------------------------------------------------------
%% Package Imports
%% \usepackage[options]{name} loads a package from the preamble.
%% Packages extend LaTeX — fonts, colors, math, figures, tables…
%% ------------------------------------------------------------
%% Latin Modern: freely scalable version of Computer Modern.
%% Prevents "Font ... not available, size <X> substituted" warnings
%% that pgfplots can trigger when computing tiny axis-label sizes.
\usepackage{lmodern}
%% Japanese font setup — change the option to match your system.
%%   Options: [noto-otf], [ipaex], [hiragino-pron], [kozuka-pr6n]
%%   No option → use the system default font
\usepackage{luatexja-preset}

%% Page layout — 'geometry' sets margins directly (top/bottom/left/right)
\usepackage[top=30mm, bottom=25mm, left=25mm, right=25mm]{geometry}

%% Color — 'svgnames' unlocks named colors like SteelBlue, Coral; 'table' colors table rows
\usepackage[svgnames,table]{xcolor}

%% Image inclusion — provides \includegraphics[width=...]{file}
\usepackage{graphicx}

%% TikZ / pgfplots — draw vector diagrams and data charts inline in LaTeX
\usepackage{tikz}
\usepackage{pgfplots}
\pgfplotsset{compat=1.18}

%% Source code — use \lstinputlisting{file} or lstlisting environment for highlighted code
\usepackage{listings}

%% Hyperlinks — makes ToC entries, \ref, \cite, and \url clickable in the PDF
\usepackage{hyperref}

%% Header and footer customization (configured in the Page Style section below)
\usepackage{fancyhdr}

%% Captions and sub-figures — sub-figures are labeled (a), (b), … automatically
%% 'format=hang' + explicit options suppress the "unknown document class" warning
%% that caption produces when it doesn't recognise ltjsarticle.
\usepackage[format=hang, font=small, labelfont=bf]{caption}
\usepackage{subcaption}

%% Mathematics — amsmath adds equation/align environments; amssymb adds ∀, ∈, ℝ, etc.
\usepackage{amsmath}
\usepackage{amssymb}

%% Professional table rules — \toprule, \midrule, \bottomrule instead of \hline
\usepackage{booktabs}

%% URL
\usepackage{url}

%% List customization — control bullet/numbered list labels, spacing, and indentation
\usepackage{enumitem}

%% ------------------------------------------------------------
%% Source Code Block Appearance
%% \definecolor{name}{model}{value} creates a named color.
%% \lstset{} applies settings globally to every code listing.
%% ------------------------------------------------------------
\definecolor{codebg}{RGB}{248,248,248}
\definecolor{codeframe}{RGB}{180,180,180}
\definecolor{kwcolor}{RGB}{0,0,180}
\definecolor{cmcolor}{RGB}{100,100,100}
\definecolor{stcolor}{RGB}{160,0,0}

\lstset{
  backgroundcolor  = \color{codebg},
  basicstyle       = \ttfamily\small,
  keywordstyle     = \color{kwcolor}\bfseries,
  commentstyle     = \color{cmcolor}\itshape,
  stringstyle      = \color{stcolor},
  numbers          = left,
  numberstyle      = \tiny\color{gray},
  numbersep        = 8pt,
  frame            = single,
  rulecolor        = \color{codeframe},
  breaklines       = true,
  breakatwhitespace= false,
  showstringspaces = false,
  tabsize          = 4,
  captionpos       = b,
  xleftmargin      = 2em,
  framexleftmargin = 1.5em,
}

%% ------------------------------------------------------------
%% Page Style (Headers & Footers)
%% \pagestyle{fancy} activates fancyhdr for the whole document.
%% \lhead/\chead/\rhead = left/center/right header slots.
%% \lfoot/\cfoot/\rfoot = left/center/right footer slots.
%% \fancypagestyle{name}{...} defines a named style applied
%% per-page with \thispagestyle{name}.
%% ------------------------------------------------------------
\pagestyle{fancy}
\fancyhf{}
\rhead{\small プログラミング実験レポート}
\lhead{\small \nouppercase{\leftmark}}
\cfoot{\thepage}
\renewcommand{\headrulewidth}{0.4pt}

%% Title page — clear all header/footer slots and suppress the horizontal rule
\fancypagestyle{titlepage}{
  \fancyhf{}
  \renewcommand{\headrulewidth}{0pt}
}

%% Table of contents pages — show only a page number, no header rule
\fancypagestyle{frontmatter}{
  \fancyhf{}
  \cfoot{\thepage}
  \renewcommand{\headrulewidth}{0pt}
}

%% ------------------------------------------------------------
%% Hyperlink Configuration
%% \hypersetup{} controls link colors and PDF metadata
%% (title, author, subject appear in PDF reader properties).
%% Load hyperref LAST — it patches many other packages.no no no i mean 
%% ------------------------------------------------------------
\hypersetup{
  colorlinks = true,
  linkcolor  = NavyBlue,
  citecolor  = NavyBlue,
  urlcolor   = NavyBlue,
  pdftitle   = {プログラミング実験レポート},
  pdfauthor  = {山田 太郎},
  pdfsubject = {Pythonを用いた気温データの統計解析},
}

%% ------------------------------------------------------------
%% Custom Commands
%% \newcommand{\name}[args]{definition} creates a reusable macro.
%% \rule{width}{thickness} draws a solid horizontal line.
%% ------------------------------------------------------------
\newcommand{\HRule}{\rule{\linewidth}{0.5mm}}

%% ============================================================
%% Document Metadata
%% These are intentionally empty — the title page is built
%% manually inside the titlepage environment below,
%% rather than using \maketitle.
%% ============================================================
\title{} % title page is built manually in the titlepage environment below
\date{}
\author{}

%% ============================================================
%% Document Body
%% Everything between \begin{document} and \end{document} is
%% the actual output. Everything above is the preamble (setup).
%% ============================================================
\begin{document}

%% ============================================================
%% Title Page
%% 'titlepage' environment creates a standalone unnumbered page.
%% \centering centres content; \vspace adds vertical space;
%% \vfill pushes the remaining content to the bottom of the page.
%% ============================================================
\begin{titlepage}
  \thispagestyle{titlepage}
  \centering
  \vspace*{2cm}

  {\large 情報工学科 プログラミング実験}\\[0.5em]
  {\large 第3回 レポート}\\[2em]

  \HRule{}\\[1em]
  {\LARGE \textbf{Pythonを用いた気温データの統計解析}}\\[0.5em]
  \HRule{}

  \vspace{3cm}

  \begin{tabular}{rl}
    \textbf{学籍番号} & ：XXXXXXXX \\[0.4em]
    \textbf{氏　　名} & ：山田 太郎 \\[0.4em]
    \textbf{所　　属} & ：情報工学科 3年 2組 \\[0.4em]
    \textbf{担当教員} & ：XX 先生 \\
  \end{tabular}

  \vfill

  \textbf{提出日：}\today

\end{titlepage}

%% ============================================================
%% Table of Contents
%% \tableofcontents auto-generates the ToC from \section/\subsection.
%% \pagenumbering{roman} switches to i, ii, iii … numbering.
%% \pagenumbering{arabic} switches back to 1, 2, 3 … and resets
%% the counter. \setcounter{page}{1} forces the count to start at 1.
%% ============================================================
\newpage
\thispagestyle{frontmatter}
\pagenumbering{roman}
\tableofcontents

\newpage
\pagenumbering{arabic}
\setcounter{page}{1}

%% ============================================================
\section{はじめに}\label{sec:intro}
%% ============================================================

\subsection{背景}

近年，気象データの解析は環境問題・農業・都市計画など多様な分野で重要性を増している．
2022年の日本における気温データには，猛暑や寒波など様々な気象現象が記録されており，
統計的な観点からの分析が求められている．

本実験では，Pythonプログラミング言語とそのライブラリを活用し，
2022年の日別平均気温データに対して統計解析とグラフ可視化を行う．

\subsection{目的}

本実験の目的を以下に列挙する．

\begin{enumerate}[label=(\arabic*)] % chktex 36
  \item Pythonを用いたテキストデータの読み込みと前処理の実装
  \item NumPyによる統計量（平均・分散・最大・最小）の計算
  \item Matplotlibを用いたデータ可視化技術の習得
  \item 実験結果に基づく考察能力の向上
\end{enumerate}

\subsection{レポートの構成}

本レポートの構成は以下の通りである．
第\ref{sec:env}節では実験環境を示す．
第\ref{sec:method}節では実装方法の詳細を述べる．
第\ref{sec:result}節では実験結果を示す．
第\ref{sec:discussion}節では結果に対する考察を行い，
第\ref{sec:summary}節でまとめを述べる．

%% ============================================================
\section{実験環境}\label{sec:env}
%% ============================================================

実験に使用したハードウェアおよびソフトウェアの環境を表~\ref{tab:env}に示す．

\begin{table}[htbp]
  \centering
  \caption{実験環境の一覧}\label{tab:env}
  \begin{tabular}{ll}
    \toprule
    \textbf{項目}    & \textbf{詳細} \\
    \midrule
    OS               & Ubuntu 22.04 LTS \\
    CPU              & Intel Core i7-1165G7 \\ % chktex 8
    メモリ           & 16 GB DDR4 \\
    Python           & 3.10.12 \\
    NumPy            & 1.24.0 \\
    Matplotlib       & 3.7.0 \\
    pandas           & 2.0.0 \\
    \bottomrule
  \end{tabular}
\end{table}

\subsection{使用データ}

使用したデータを表~\ref{tab:data}に示す．

\begin{table}[htbp]
  \centering
  \caption{使用データの概要}\label{tab:data}
  \begin{tabular}{lll}
    \toprule
    \textbf{ファイル名}       & \textbf{形式}   & \textbf{内容} \\
    \midrule
    \texttt{Temperature2022.txt} & テキスト (UTF-8) & 2022年 日別平均気温 \\
    \bottomrule
  \end{tabular}
\end{table}

%% ============================================================
\section{実験方法・実装}\label{sec:method}
%% ============================================================

\subsection{プログラムの全体構成}

実装したプログラムは3つの関数で構成される。
図~\ref{fig:flowchart}にプログラムの処理フローを示す。

\begin{figure}[htbp]
  \centering
  \begin{tikzpicture}[
    node distance   = 1.4cm,
    box/.style      = {rectangle, draw=SteelBlue, fill=SteelBlue!10,
                       minimum width=5.5cm, minimum height=0.8cm,
                       align=center, text centered, font=\small},
    decision/.style = {diamond, draw=DarkOrange, fill=Orange!10,
                       aspect=2.5, text centered, font=\small},
    arrow/.style    = {->, >=stealth, thick, draw=gray!70},
  ]
    \node[box] (start)  {開始};
    \node[box, below of=start] (load) {\texttt{load\_temperature()} \\ ファイル読み込み・前処理};
    \node[box, below of=load]  (stat) {\texttt{calc\_statistics()} \\ 統計量の計算};
    \node[box, below of=stat]  (plot) {\texttt{plot\_temperature()} \\ グラフ生成・保存};
    \node[box, below of=plot]  (end)  {終了};

    \draw[arrow] (start) -- (load);
    \draw[arrow] (load)  -- (stat);
    \draw[arrow] (stat)  -- (plot);
    \draw[arrow] (plot)  -- (end);
  \end{tikzpicture}
  \caption{プログラムの処理フロー}\label{fig:flowchart}
\end{figure}

\subsection{ソースコードの説明}

リスト~\ref{lst:analysis}に実装したPythonプログラムを示す．
ファイル \texttt{code/analysis.py} として保存されており，
\texttt{python code/analysis.py} で実行できる．

%% 外部ファイルの読み込みによるコード挿入
\lstinputlisting[
  language = Python,
  caption  = {気温データ解析プログラム（\texttt{code/analysis.py}）},
  label    = lst:analysis
]{code/analysis.py}

\subsection{統計量の計算式}

NumPyを活用することで，以下の統計量を計算する．

\textbf{平均値}（算術平均）：
\begin{equation}
  \bar{x} = \frac{1}{n} \sum_{i=1}^{n} x_i
  \label{eq:mean}
\end{equation}

\textbf{標準偏差}：
\begin{equation}
  \sigma = \sqrt{\frac{1}{n} \sum_{i=1}^{n} {(x_i - \bar{x})}^{2}}
  \label{eq:std}
\end{equation}

ここで，$n$ はデータ数，$x_i$ は $i$ 日目の気温（℃）を表す．

%% ============================================================
\section{実験結果}\label{sec:result}
%% ============================================================

\subsection{統計量の結果}

プログラムを実行して得られた統計量を表~\ref{tab:stats}に示す．

\begin{table}[htbp]
  \centering
  \caption{2022年 気温データの統計量}\label{tab:stats}
  \rowcolors{2}{gray!10}{white}
  \begin{tabular}{lr}
    \toprule
    \textbf{統計量} & \textbf{値 (℃)} \\
    \midrule
    平均値          & 15.23 \\
    標準偏差        & 8.47  \\
    最大値          & 35.6  \\
    最小値          & $-3.2$ \\
    中央値          & 16.1  \\
    第1四分位数     & 8.4   \\
    第3四分位数     & 23.1  \\
    \bottomrule
  \end{tabular}
\end{table}

\subsection{月別平均気温}

図~\ref{fig:monthly}に月別平均気温のグラフを示す．

\begin{figure}[htbp]
  \centering
  \begin{tikzpicture}
    \begin{axis}[
      width         = 0.88\linewidth,
      height        = 6.2cm,
      title         = {2022年 月別平均気温},
      xlabel        = {月},
      ylabel        = {気温 (℃)},
      xtick         = {1,2,3,4,5,6,7,8,9,10,11,12},
      xticklabels   = {1月,2月,3月,4月,5月,6月,7月,8月,9月,10月,11月,12月},
      xticklabel style = {rotate=45, anchor=east, font=\small},
      ymin          = -5, ymax = 35,
      grid          = both,
      grid style    = {line width=0.2pt, draw=gray!30},
      major grid style = {line width=0.4pt, draw=gray!60},
      legend pos    = north west,
      legend style  = {font=\small},
      tick label style = {font=\small},
    ]
      \addplot[color=SteelBlue, mark=*, thick, mark size=2.5pt] coordinates {
        (1, 3.2)  (2, 5.1)  (3, 9.4)  (4, 14.7) (5, 19.2) (6, 23.1)
        (7, 28.5) (8, 29.8) (9, 24.3) (10, 17.6) (11, 11.2) (12, 5.8)
      };
      \addlegendentry{月別平均気温}

      %% Annual average reference line (dashed horizontal)
      \addplot[color=Crimson, dashed, thick, domain=1:12] {15.23};
      \addlegendentry{年平均 (15.23℃)}
    \end{axis}
  \end{tikzpicture}
  \caption{2022年 月別平均気温の変化}\label{fig:monthly}
\end{figure}

\subsection{気温の分布}

図~\ref{fig:hist}にプログラムが生成した気温の度数分布（ヒストグラム）を示す．

\begin{figure}[htbp]
  \centering
  \begin{tikzpicture}
    \begin{axis}[
      width         = 0.75\linewidth,
      height        = 5cm,
      title         = {気温の度数分布},
      xlabel        = {気温 (℃)},
      ylabel        = {日数},
      ybar,
      bar width     = 5pt,
      grid          = both,
      grid style    = {line width=0.2pt, draw=gray!30},
      ymajorgrids   = true,
      tick label style = {font=\small},
    ]
      \addplot[fill=Coral!70, draw=DarkRed!60, opacity=0.85] coordinates {
        (-5,2)  (-3,5)  (-1,8)  (1,12)  (3,15)  (5,18)  (7,20)
        (9,22)  (11,25) (13,28) (15,30) (17,27) (19,24) (21,20)
        (23,18) (25,15) (27,12) (29,9)  (31,6)  (33,3)  (35,1)
      };
    \end{axis}
  \end{tikzpicture}
  \caption{気温の度数分布（ヒストグラム）}\label{fig:hist}
\end{figure}

%% --- Example: insert an external image from the images/ folder ---
%% \includegraphics[width=0.82\linewidth]{file} scales the image
%% to 82% of the text width. The 'figure' float places it automatically.
%% \begin{figure}[h]
%%   \centering
%%   \includegraphics[width=0.82\linewidth]{images/temperature.png}
%%   \caption{Matplotlibで生成した気温グラフ}
%%   \label{fig:matplotlib}
%% \end{figure}

%% --- Example: two side-by-side sub-figures using subcaption ---
%% \hfill inserts flexible horizontal space to push figures apart.
%% \begin{figure}[h]
%%   \centering
%%   \begin{subfigure}[b]{0.48\linewidth}
%%     \includegraphics[width=\linewidth]{images/line.png}
%%     \caption{折れ線グラフ}
%%   \end{subfigure}
%%   \hfill
%%   \begin{subfigure}[b]{0.48\linewidth}
%%     \includegraphics[width=\linewidth]{images/histogram.png}
%%     \caption{ヒストグラム}
%%   \end{subfigure}
%%   \caption{気温データの可視化結果}
%%   \label{fig:sub}
%% \end{figure}

%% ============================================================
\section{考察}\label{sec:discussion}
%% ============================================================

\subsection{統計量の解釈}

表~\ref{tab:stats}の結果から，以下のことが考察される．

\begin{itemize}
  \item 年間平均気温は 15.23℃ であり，日本の年平均気温（約15℃）とほぼ一致した~\cite{jma}．
  \item 標準偏差が 8.47℃ と大きく，季節変動が著しいことを表している．
        式~(\ref{eq:std})で示すように，標準偏差は各データが平均値からどれだけ
        ばらついているかを表す指標である．
  \item 最大値 35.6℃，最小値 −3.2℃ の差は 38.8℃ に及び，
        年間を通じた気温の寒暖差が大きいことが確認できた．
  \item 図~\ref{fig:hist}の分布形状は正規分布に近い形を示しており，
        気温データが統計的な解析手法に適していることが示唆される．
\end{itemize}

\subsection{プログラムの評価}

実装したプログラム（リスト~\ref{lst:analysis}）について，以下の観点から評価を行った．

\begin{description}
  \item[正確性]
        NumPyの組み込み関数（\texttt{np.mean}, \texttt{np.std} 等）を使用したため，
        浮動小数点誤差を最小限に抑えた数値計算が実現できた．

  \item[可読性]
        関数に処理を分割した構造化プログラミングを採用しており，
        各関数の責務が明確で保守性が高い．

  \item[堅牢性]
        ファイルが存在しない場合にサンプルデータを生成するフォールバック処理を
        組み込み，エラー耐性を高めた．

  \item[改善点]
        現実装では \texttt{for} ループによるデータ読み込みを行っているため，
        大規模データに対する処理速度が課題となる．
        pandasの \texttt{read\_csv} や NumPyの \texttt{loadtxt} への移行で
        高速化が期待される．
\end{description}

\subsection{今後の発展}

本実験の成果を基に，以下の発展的な取り組みが考えられる．

\begin{enumerate}[label=\arabic*.]
  \item 複数年度（2019〜2022年）のデータを用いた経年変化の分析
  \item 機械学習（線形回帰・LSTM等）を用いた気温予測モデルの構築
  \item 気温以外の気象要素（降水量・日照時間）との相関分析
\end{enumerate}

%% ============================================================
\section{まとめ}\label{sec:summary}
%% ============================================================

本実験では，Pythonを用いて2022年の気温データに対する統計解析プログラムを実装し，
実験・考察を行った．得られた成果を以下にまとめる．

\begin{enumerate}[label=(\arabic*)] % chktex 36
  \item テキストファイルの読み込みと前処理を実装し，
        365日分の日別平均気温データを正常に取得できた．
  \item NumPyを用いた統計計算により，年間の気温特性（平均 15.23℃，
        標準偏差 8.47℃）を定量的に把握した．
  \item Matplotlibによる折れ線グラフおよびヒストグラムの生成に成功し，
        データの視覚的な分析が可能となった．
  \item 考察を通じて，データの分布特性や季節変動パターンを明らかにした．
\end{enumerate}

今後の課題として，より大規模なデータセットへの対応と，
機械学習を応用した予測モデルの構築に取り組む予定である．

%% ============================================================
%% Bibliography
%% 'thebibliography' is the manual citation environment.
%% \bibitem{key} defines an entry referenced in text by \cite{key}.
%% For bigger projects, consider BibTeX (.bib files) or BibLaTeX.
%% ============================================================
\begin{thebibliography}{9}

\bibitem{python}
  Python Software Foundation,
  \textit{Python 3.10 Documentation},
  \url{https://docs.python.org/3/},
  2022.

\bibitem{numpy}
  C. R. Harris \textit{et al.},
  ``Array programming with NumPy'',
  \textit{Nature}, vol.~585, pp.~357--362, 2020.

\bibitem{matplotlib}
  J. D. Hunter,
  ``Matplotlib: A 2D Graphics Environment'',
  \textit{Computing in Science \& Engineering},
  vol.~9, no.~3, pp.~90--95, 2007.

\bibitem{pandas}
  The pandas development team,
  \textit{pandas-dev/pandas: Pandas},
  Zenodo, 2023.
  \url{https://doi.org/10.5281/zenodo.3509134}

\bibitem{jma}
  気象庁,
  ``過去の気象データ・ダウンロード'',
  \url{https://www.data.jma.go.jp/gmd/risk/obsdl/},
  参照日：\today.

\end{thebibliography}

\end{document}
